\section{Reproducibility and artifact availability}
\label{sec:reproducibility}

% Figure 9 -- the integrity chain, from freeze to published numbers. No infrastructure addresses or account
% names appear; hosts are referred to by role only.
\begin{figure}[t]
\centering
\begin{tikzpicture}[font=\footnotesize, x=1mm, y=1mm,
  stg/.style={blk, minimum width=26mm, minimum height=12mm, align=center},
  chk/.style={blkacc, minimum width=26mm, minimum height=9mm, align=center, font=\scriptsize}]

\node[stg] (freeze) at (16,34) {\textbf{Freeze}\\[-1pt]\scriptsize mechanism, runner,\\[-2pt]\scriptsize analysis};
\node[stg] (seal)   at (50,34) {\textbf{Seal}\\[-1pt]\scriptsize benchmark hashed,\\[-2pt]\scriptsize gold fixed};
\node[stg] (prereg) at (84,34) {\textbf{Pre-register}\\[-1pt]\scriptsize hypotheses, test,\\[-2pt]\scriptsize corrections};
\node[stg] (run)    at (118,34){\textbf{Run}\\[-1pt]\scriptsize 18 cells,\\[-2pt]\scriptsize 57{,}888 runs};
\node[stg] (anal)   at (150,34){\textbf{Analyse}\\[-1pt]\scriptsize frozen script,\\[-2pt]\scriptsize 100{,}000 permutations};

\foreach \a/\b in {freeze/seal, seal/prereg, prereg/run, run/anal} \draw[flow] (\a) -- (\b);

\node[chk] (c1) at (16,17) {file hashes match\\\cond{FREEZE} manifest};
\node[chk] (c2) at (50,17) {\cond{sha256sum -c}\\on the sealed set};
\node[chk] (c3) at (84,17) {document hash\\matches the freeze};
\node[chk, minimum height=12mm] (c4) at (118,17)
  {gates G1--G4: all cells\\done, exact row counts,\\clean exits, run identity};
\node[chk] (c5) at (150,17) {independent\\cross-check agrees};

\foreach \s/\c in {freeze/c1, seal/c2, prereg/c3, run/c4, anal/c5} \draw[flowacc] (\s) -- (\c);

\node[lbl, text=cEBG, anchor=north, align=center] at (50,8)
  {\scriptsize checked at \textbf{every one} of the 20 job starts, not once at the end};

\node[blkwarn, minimum width=54mm, minimum height=13mm, align=center] (recov) at (118,2)
  {\scriptsize \textbf{Recovery, recorded not hidden:} a node reboot interrupted\\[-1pt]
   \scriptsize two cells. The damaged tail of each output file was\\[-1pt]
   \scriptsize quarantined (0 complete rows lost) and the cells resumed.\\[-1pt]
   \scriptsize A reporting-only check (G5) records the platform transition.};
\draw[flowwarn] (c4.south) -- (118,8.5);

\node[blkplain, minimum width=44mm, minimum height=8mm, align=center] (man) at (140,-14)
  {\scriptsize \textbf{SHA-256 manifest} over every\\[-1pt]\scriptsize input, output, script and report};
\draw[flow] (c5.south) -- (152,-8) -- (140,-8) -- (140,-10);
\end{tikzpicture}

\vspace{1mm}
\caption[Experimental integrity chain]{%
\textbf{The integrity chain.} Each stage has a check that must pass before the next begins, and the checks on the
first three stages were re-run at every one of the 20 job starts rather than once at the end: all 20 reported zero
mechanism-file mismatches and the same pre-registration hash. The gates before analysis verify that all 18 cells
finished, that each holds exactly 268 $\times$ 12 unique (task, condition) pairs with no duplicate and no
unparsable line, that every cell's last job exited cleanly, and that the recorded run identity --- model, snapshot,
decoding, seed, dtype, parser, framework versions, and the hashes of all 14 mechanism files --- is identical across
every start of a cell. One deviation is recorded rather than smoothed over: an unplanned reboot interrupted two
cells, which resumed after the damaged tail of each file was quarantined, and a separate reporting-only check
documents the resulting platform transition (\Cref{sec:limitations}). Host names and addresses are deliberately
absent from this figure and from the released artifacts.\\
\emph{Accessible description:} a five-stage left-to-right chain, each stage with a check box beneath it, plus a
recovery note under the run stage and a final hash manifest.\\
\emph{Source:} \srcfile{PHASE3/scripts/finalize_gates.py}, \texttt{verify\_phase3.py},
\texttt{platform\_continuity.py}; \srcfile{PHASE3/PHASE3_SHA256_MANIFEST.txt}.}
\label{fig:integrity}
\end{figure}


\subsection{Provenance of every row}

Each cell writes a manifest at every start recording the model identity and resolved snapshot, the adapter (none
here) and its hash, decoding parameters and seed, dtype, maximum new tokens, parser version, family adapter,
framework and CUDA runtime versions, the GPU model, the split hash, the condition list, the SHA-256 of all 14
mechanism and harness files, the job id, the start time, and how many pairs were already complete at that start.
The 14 code hashes are identical across all 20 starts of the campaign. A SHA-256 manifest over every raw file is
written when the campaign closes and re-verified before anything is analysed (\Cref{fig:integrity}).

\subsection{Regenerating every number}

The whole pipeline runs as one CPU-only batch job with no interactive steps: integrity gates, staging, the
platform-continuity report, verification, the frozen analysis, the exploratory analyses, an independent
cross-check, report generation and a final hash manifest, stopping at the first failure. Every reported number
comes from that path; the independent cross-check must agree exactly or the pipeline stops, and the report
generator computes no statistic of its own. On the shareable development split the same pipeline reproduces its
outputs byte for byte.

\subsection{What a reproduction can and cannot match}

The \emph{analysis} is reproducible byte for byte from the archived raw outputs: permutations and bootstraps are
seeded and the gates re-verify the inputs first. Regenerating the \emph{raw outputs} is a weaker guarantee, with
three specific caveats: GPU kernels are not bitwise deterministic across hardware or library versions; two cells
were not produced on a single host platform and cannot be reproduced row for row on one driver version; and one of
those cells re-seeded its sampler when it resumed, so reproducing it byte for byte would require reproducing the
interruption at the same row (\Cref{tab:deviations}).

\subsection{Availability}

\textbf{Released with the paper:} the code (mechanism, harness, runner, analysis, integrity gates, tests), the
frozen manifests, the pre-registration, the development splits and their runs, all aggregate results in JSON and
CSV, the generated reports, the deviation ledger and the platform-continuity report, and a reproducibility package
documenting environment, model snapshots, dataset rules, batch scripts, analysis instructions, expected outputs,
tests, limitations and deviations, with SHA-256 manifests over all of it.

\textbf{Withheld until publication:} the eval3 items, their annotations and the per-run transcripts, so that the
test set stays uncontaminated; their hashes are public, so a later release can be checked against the frozen
versions. The cluster's per-run logs are withheld for the same reason, because their progress lines carry per-task
outcomes.

\textbf{Not redistributed:} model weights. The exact snapshot hashes are listed in the package; obtain the weights
under each model's own licence. Because the eval3 texts were written by Qwen2.5-14B-Instruct, that model's licence
terms on outputs apply to them.

Nothing here claims public availability for an artifact that is in fact restricted, and the released materials
contain no credentials, tokens, keys, host names, addresses or account identifiers; a check in the packaging
script fails the build if any appear.
